\documentclass[11pt]{article}
\usepackage[margin=0.85in]{geometry}
\usepackage{times}
\usepackage{parskip}
\setlength{\parskip}{4pt}
\usepackage[hidelinks]{hyperref}
\pagestyle{empty}

\begin{document}

\noindent\textbf{Cover Letter}\\[2pt]
\noindent To: Guest Editors, IEEE Journal of Selected Topics in Signal Processing\\
\noindent Special Issue: \emph{Autonomous and Evolutive Optimization in Networked AI}\\
\noindent (Lead Guest Editor: Prof.\ Liang Song; with Profs.\ J.\ Liu, A.\ Dvir,
A.\ Skodras, V.\ C.\ M.\ Leung, and Q.\ Bi)\\[6pt]

\noindent Dear Guest Editors,

We are pleased to submit our manuscript, \textbf{``Interpretable Federated
Kolmogorov--Arnold Surrogates for Autonomous and Evolutive Resource Optimization
over Non-Stationary LEO Links,''} for consideration in the Special Issue on
\emph{Autonomous and Evolutive Optimization in Networked AI}.

\textbf{Scope and fit.} The Special Issue calls for online, self-evolving
mechanisms that let distributed models improve in multi-agent, time-varying
environments, with emphasis on the foundations of signal processing for networked
AI, networked AI for cognitive communications, and online drift detection and
compensation. Our paper addresses exactly this setting. We study autonomous
resource optimization over low Earth orbit (LEO) links, where the per-instance
optimum has no closed form and the operating regime drifts as satellites pass and
interference handovers occur. We cast the problem as learning the \emph{solution
map} of an energy-efficiency power-allocation problem, harden it into a
mixed-integer program, and decompose it into an inner statistical-learning task and
an outer communication-control task whose static optimum we derive in closed form.
A federated Kolmogorov--Arnold network learns the solution map online, while a
contextual-bandit controller \emph{learns the Lagrangian dual price} of the uplink
budget, unifying supervised learning with online reinforcement, as the call
envisions.

\textbf{Contributions and findings.} Over twenty seeds with paired significance
tests, the KAN surrogate lowers prediction error by about $1.8\times$ and
energy-efficiency regret by about $2\times$ versus a parameter-matched multilayer
perceptron (Wilcoxon $p<0.01$), while a linear predictor fails entirely, confirming
the map is strongly nonlinear. A learned sparsification controller halves the
uplink volume while still outperforming the perceptron, and the KAN extrapolates to
unseen interference regimes roughly an order of magnitude more accurately. Uniquely,
the additive-spline structure yields closed-form, auditable design rules that a
black-box network cannot provide. We compare against a model-based deep-unfolding
learn-to-optimize baseline and analyze each component through one-factor ablations.

\textbf{Rigor and transparency.} All results are multi-seed with standard
deviations and Wilcoxon tests, every comparison is parameter-matched, and we report
transparently the cases where our method does \emph{not} help (online grid
self-evolution gives no significant gain; under abundant data or large CSI error
the advantage narrows). All code and result files that reproduce every number and
figure are openly available.

\textbf{Declarations.} This manuscript is original, is not under consideration
elsewhere, all authors approve the submission, and we declare no conflicts of
interest. The work was supported in part by PTIT. We believe the paper is a strong
fit for the Special Issue and thank you for your consideration.

\vspace{8pt}
\noindent Sincerely,\\[4pt]
\noindent Phuc Hao Do (corresponding author)\\
\noindent Center for Advanced Innovation, Research and Application,\\
\noindent Da Nang Architecture University, Da Nang, Vietnam\\
\noindent E-mail: haodp@dau.edu.vn\\[4pt]
\noindent on behalf of co-authors Huu Phu Le and Truong Duy Dinh

\end{document}
